\documentclass[10pt,a4paper]{amsart}
\usepackage[margin=19mm]{geometry}
\usepackage{amsmath,amssymb,mathtools}
\usepackage[scr=rsfs,cal=euler]{mathalfa}
\usepackage{xfrac}
\usepackage[unicode,hidelinks,bookmarks=false]{hyperref}
\newcommand{\ie}{i.e.\ }
\newcommand{\cf}{cf.\ }
\newcommand{\resp}{respectively\ }
\newcommand{\Subsection}{\subsection}
\providecommand{\nobreakdash}{\nobreak\hbox{-}}
\setlength{\emergencystretch}{2em}
\multlinegap0pt
\usepackage{bm}
\usepackage{bbm}
\usepackage{dsfont}
\usepackage{xspace}
\usepackage{cancel}
\usepackage{mathtools}
\usepackage{amsthm}
\makeatletter
\@ifundefined{thm}{\newtheorem{thm}{Theorem}}{}
\makeatother
\title[Real non-attractive fixed point conjecture for complex harmo]{Real non-attractive fixed point conjecture for complex harmonic functions}
\author{Mohd Vaseem}
\date{Source: arXiv:2507.18414v1 (CC BY 4.0)}
\begin{document}
\maketitle

\noindent\emph{Source and licence.} Real non-attractive fixed point conjecture for complex harmonic functions. Version arXiv:2507.18414v1 (CC BY 4.0), distributed under the Creative Commons Attribution 4.0 International licence, as recorded in the arXiv licence metadata for this exact version and shown on the abstract page https://arxiv.org/abs/2507.18414v1. Full rights notice: licenses/arxiv-2507.18414-CC-BY-4.0.txt.\par\smallskip

\noindent\textit{Editorial scope.} Counted unit: the Thm whose source label is \texttt{T1} in the article (source0.tex lines 223--228, source label \texttt{T1}), delivered together with the article's own complete original proof of it (source0.tex lines 230--267, one \texttt{proof} environment of 38 lines). No mathematics is written, repaired or re-derived: statement and proof are the article's own text line for line. Required context retained uncounted: none. Every definition, display and label of those lines is retained unchanged. Omitted from this package and not redistributed: 1--222, 229--229, 268--415. Nothing inside a retained excerpt is omitted or altered: every retained body is delivered exactly as the article's lines are written, so no omission receipt of any kind accompanies this package. Citations: the retained excerpts carry no citation command at all, so no bibliography entry of the article is transcribed and no reference is redistributed. Reference adaptation: none was needed and none was made; every reference inside the delivered unit resolves inside it, so no unresolved reference or citation is produced. Printed slips kept literal: no printed word, symbol, hypothesis or number of the counted statement or of its proof was altered. Layout and class: the article's own class is \texttt{amsart} with packages including amsmath, amsthm, amssymb, hyperref, cleveref, mathrsfs, graphicx, comment; the wrapper is the lane's pinned \texttt{amsart} editorial wrapper at 10pt on A4, which restates the theorem-like environments the retained text needs and adds the article's own macro definitions that the retained text needs (none), with extra wrapper packages (bm, bbm, dsfont, xspace, cancel, mathtools). The delivered numbering is the unit's own. Assets: no retained span contains a figure environment, a graphics-inclusion command, a PDF-page inclusion, a table or a TikZ picture; no image file is copied into this package, the package declares no asset dependency and no image file is redistributed with it. Licence: version arXiv:2507.18414v1 of this article is distributed under the Creative Commons Attribution 4.0 International licence, as recorded in the arXiv licence metadata for this exact version and shown on the abstract page https://arxiv.org/abs/2507.18414v1. A full rights notice is delivered at licenses/arxiv-2507.18414-CC-BY-4.0.txt. Characters: every retained excerpt is pure ASCII, so the delivered engine, XeLaTeX with its default Latin Modern text font, reports no missing character.
\begin{thm}\label{T1}


Every rational harmonic function $f=h+\overline{g}$ where both the rational functions $h$ and $g$ are of degree greater than or equal to two and with a super-attracting fixed point, has a $\mathfrak{h}$-fixed point  $\zeta=\mu+\overline{\omega}$ such that the real part of its multipliers is greater than or equal to $1$, i.e., $Re\left( \lambda=\partial_{z}h(\mu)\right)\geq 1$ and $Re\left( \theta=\partial_{z}g(\mu)\right)\geq 1$. In particular, this is true for all polynomial harmonic functions with degree greater than or equal to two.
    
\end{thm}

\begin{proof}

Let $f=h+\overline{g}$ be rational harmonic function, where both the rational functions, $h$ and $g$ are of degree greater than or equal to two and with a super-attracting fixed point. 



Now firstly we are trying to show for the rational function $h$ that there exist a fixed point whose real part is greater than equal to one. If a rational function $h$ has a multiple fixed point then in this case it's multiplier is $1$ and real part of its multiplier is one.\\

Now consider the case when the rational function $h$ has all of its fixed points simple (Multipliers different from one). If rational function $h$ is of degree $d$ then it has $d+1$ distinct fixed points. Suppose that the multiplier $\lambda_{i}=\partial_{z}h(\mu_{i})$ for $i=1,2,3,\ldots,d+1$, where $h(\mu_{i})=\mu_{i}$ for $i=1,2,3,\ldots, d+1.$ 

Now by  using the rational fixed point theorem we have, $\displaystyle\sum_{i=1}^{d+1}\frac{1}{1-\lambda_{i}}=1$. However, if the rational function $h$ has a super attracting fixed point, say $\mu_{d+1}$, for this point its multiplier is $\lambda_{d+1}=0$. Therefore,
\begin{equation}\label{E3}
\displaystyle\sum_{i=1}^{d}\frac{1}{1-\lambda_{i}}=0
\end{equation}
Comparing the real part, we get that
\begin{align*}
Re\left( \frac{1}{1-\lambda_{i}}\right)&=\frac{1}{2}\left[ \frac{1}{1-\lambda_{i}}+\frac{1}{1-\bar{\lambda_{i}}}\right]\quad\quad\quad(\lambda_{i}=x+iy) \\
&=\frac{1-\bar{\lambda_{i}}+1-\lambda_{i}}{2(1-\lambda_{i})\overline{(1-\lambda_{i})}}\\
&=\frac{1-x+iy+1-x-iy}{2{\vert 1-\lambda_{i}\vert}^{2}}\\
&=\frac{1-x}{{\vert 1-\lambda_{i}\vert}^{2}}\\
&=\frac{Re(1-\bar{\lambda_{i}})}{{\vert 1-\lambda_{i}\vert}^{2}}
\end{align*}
Hence by equation \ref{E3}

\begin{equation}\label{E4}
\displaystyle\sum_{i=1}^{d}\frac{Re(1-\bar{\lambda_{i}})}{\mid 1-\lambda_{i}
\mid^{2}}=0 
\end{equation}
Since $Re(1-\bar{\lambda_{i}})=1-Re(\lambda_{i})$ and $\mid 1-\lambda_{i}\mid^{2}>0$ for each $i$, it is not possible to have $Re(\lambda_{i})<1$ for every $i$. Therefore, there exists $j\in\lbrace1,2,\cdots,d\rbrace$ such that $Re\left( \lambda_{i}=\partial_{z}h(\mu)\right)\geq 1$.

Therefore every rational function $h$ having degree at least two and with a super attracting fixed point has a fixed point, the real part of whose multiplier is greater than or equal to $1$.

Similarly we can show the existence of a fixed of rational function $g$ whose real part of multiplier is greater than or equal to one, $Re\left( \theta_{i}=\partial_{z}g(\mu_{i})\right)\geq 1$.


 
 Therefore, both the rational functions $h$ and $g$ have fixed points, say $\mu$ and $\omega$ such that $h(\mu)=\mu$ and $g(\omega)=\omega$. Then the real part of the multipliers  $Re\left( \lambda=\partial_{z}h(\mu)\right)\geq 1$ and $Re\left( \theta=\partial_{z}g(\mu)\right)\geq 1$. Also as $\mu+\overline{\omega}=h(\mu)+\overline{g(\omega)}$. Thus $\mu+\overline{\omega}$ is a $\mathfrak{h}$- fixed point of $f$. Therefore,   Every rational harmonic function $f=h+\overline{g}$, where both the rational functions $h$ and $g$ are of degree greater than or equal to two and with a super-attracting fixed point, has a $\mathfrak{h}$-fixed point  $\zeta=\mu+\overline{\omega}$ such that the real part of its multipliers is greater than or equal to $1$.
\end{proof}

\end{document}
